\section{完成率与代理质量由不同端点领先}
\label{sec:evaluation}

\textbf{实验协议。} 三个云端端点在同一个智能体修订版和 Track-A
契约~\cite{fpt26tracka} 下处理相同的 \CorpusTaskCount{} 个均衡任务套件。
我们报告公开门控完成率，以及提交侧综合产生的 starter 锚定 QoR 代理指标。
该指标仅覆盖25个显式优化任务，并使用公式~\eqref{eq:score}；它不包含评测侧
隐藏验证和 reference 回退。

\begin{table}[!h]
  \centering
  \caption{\textbf{完成率与代理质量由不同端点领先。} QoR 代理指标为
  $100\,Q_{\mathrm{HW}}\,E$，使用25个优化任务的提交侧综合结果与 starter
  进行比较。}
  \label{tab:headline-results}
  \scriptsize
  \resizebox{\columnwidth}{!}{%
  \begin{tabular}{lrrrrr}
    \toprule
    端点 & 完成数 & 令牌（M） & 积分 & QoR 代理均值 & 代理 $>$76 \\
    \midrule
    \MainModelResultsRows
    \bottomrule
  \end{tabular}}
\end{table}

\textbf{结果。} 表~\ref{tab:headline-results} 和逐类别分解（附录
~\ref{app:results}）揭示三种模式。

（1）\textbf{Qwen3.6 的公开门控完成率最高。} Qwen3.6-27B 取得了
\QwenThreeSixSuccessRate{} 的完成率（\QwenThreeSixSuccessCount{}/
\CorpusTaskCount{} 个任务），包括结构 \CoSim{} 修复的25/25。Qwen3.5 的结构
修复完成数为20/25，是本次比较中的最低分类结果。

（2）\textbf{DeepSeek 以更高令牌成本取得最高优化代理指标。} DeepSeek 的平均
$Q_{\mathrm{HW}}$ 为 \DeepSeekQorQhw{}，代理均值为 \DeepSeekQorProxy{}，其中
\DeepSeekProxyAboveRate{} 的任务高于76。Qwen3.6 的对应结果为
\QwenThreeSixQorQhw{} 和 \QwenThreeSixQorProxy{}；Qwen3.5 为
\QwenThreeFiveQorQhw{} 和 \QwenThreeFiveQorProxy{}。DeepSeek 使用的令牌是
Qwen3.6 的3.1$\times$，Qwen3.5 使用的令牌最少。

（3）\textbf{三个端点呈现完成率、代理质量和成本之间的权衡。} Qwen3.6 的完成率
最高，DeepSeek 的优化代理指标最高，Qwen3.5 的令牌用量最低。每个端点只有一次
运行，而且托管服务商不同，因此不能把差异单独归因于模型架构。

\takeaway{Qwen3.6 的公开门控完成率最高，DeepSeek 的 starter 锚定优化代理指标
最高，Qwen3.5 使用的令牌最少。}
